\section{Deployment}
\label{sec:deployment}

The teaching assistant is delivered as a Gradio web application
(\texttt{app.py})~\cite{abid2019gradio} with four tabs, each backed by one of the four task
formats the model was trained on. The application accepts either the custom transformer or a
pretrained model selected at launch time via a CLI flag, which made A/B comparison during
development straightforward. Figure~\ref{fig:app} shows the structure.

\begin{figure}[h]
\centering
\begin{tikzpicture}[node distance=0.4cm and 0.7cm]

    \node[block, fill=green!10] (app) {Gradio app (\texttt{app.py})};

    \node[data, below=1.0cm of app, xshift=-4.6cm] (ask)   {\textbf{Ask} \\ \footnotesize free-form Q\&A};
    \node[data, below=1.0cm of app, xshift=-1.5cm] (quiz)  {\textbf{Quiz} \\ \footnotesize multiple choice};
    \node[data, below=1.0cm of app, xshift=1.5cm]  (prac)  {\textbf{Practice} \\ \footnotesize long-answer + score};
    \node[data, below=1.0cm of app, xshift=4.6cm]  (flash) {\textbf{Flashcards} \\ \footnotesize term $\rightarrow$ definition};

    \node[block, below=2.0cm of quiz, xshift=-1.6cm] (rag) {RAG retriever};
    \node[model, below=2.0cm of prac, xshift=1.6cm] (gen) {Generator \\ \footnotesize (custom / Qwen)};

    \node[data, below=0.6cm of rag, xshift=0cm] (idx2) {\texttt{rag\_index/}};
    % \node[data, below=0.6cm of gen, xshift=0cm] (prog) {\texttt{progress.json} \\ \footnotesize (--track)};

    \draw[arrow] (app) -- (ask);
    \draw[arrow] (app) -- (quiz);
    \draw[arrow] (app) -- (prac);
    \draw[arrow] (app) -- (flash);

    \draw[arrow] (ask)   -- (rag);
    \draw[arrow] (quiz)  -- (rag);
    \draw[arrow] (prac)  -- (rag);
    \draw[arrow] (flash) -- (rag);

    \draw[arrow] (ask)   -- (gen);
    \draw[arrow] (quiz)  -- (gen);
    \draw[arrow] (prac)  -- (gen);
    \draw[arrow] (flash) -- (gen);

    \draw[arrow] (rag) -- (idx2);
    % \draw[arrow] (gen) -- (prog);

\end{tikzpicture}
\caption{Deployment topology. Each tab calls the same retriever and the same generator; the
only thing that changes between tabs is the task token prepended to the prompt.}
\label{fig:app}
\end{figure}

\paragraph{Runtime modes.}
The app supports two CLI configurations. The first (\texttt{--pretrained qwen-3b}) loads
Qwen2.5-3B-Instruct as the generator without retrieval. The second adds
\texttt{--index ../rag/rag\_index} to enable RAG-backed retrieval.

\paragraph{Why Gradio.}
Gradio was chosen over a hand-built web stack because the project deliverable is the model,
not the UI. Its tab/component primitives map almost one-to-one onto the four task formats,
and a single function call serves the application from a public URL. The
retriever and generator are loaded once at startup and reused across requests, so the only
per-request cost is the embedding of the query and the autoregressive decode.

\paragraph{Public access.}
The application is exposed to the internet via a Cloudflare tunnel running on the same
machine as the GPU. At the time of writing the live URL is
\url{https://amy-appear-predict-shadow.trycloudflare.com/}.
Cloudflare tunnel addresses are ephemeral; if the address changes the current URL will be
updated in the project repository at
\url{https://github.com/Hagestregen/dat255-teaching-assistant}.

\paragraph{Maintenance.}
Because the RAG index, the generator, and the UI are decoupled, each can be replaced
independently. Re-running \texttt{chunker.py} and \texttt{embedder.py} on a new course edition
yields a fresh \texttt{rag\_index/} folder that the app picks up at next startup with no code
change. Replacing the generator is similarly local: the app accepts any model exposing the
shared \texttt{generate(prompt, task)} interface, so the comparison runs from
Section~\ref{sec:evaluation} can be reproduced live by re-launching with a different
\texttt{--pretrained} flag.
